\documentclass[10pt]{article}
\usepackage[margin=0.55in]{geometry}
\usepackage{amsmath,amssymb,booktabs,graphicx,float,microtype}
\usepackage[T1]{fontenc}
\setlength{\parindent}{0pt}
\setlength{\parskip}{3pt}
\newcommand{\ket}[1]{\lvert #1\rangle}
\newcommand{\bra}[1]{\langle #1\rvert}
\begin{document}

\begin{center}
{\Large\bfseries Experiment 09: cross-target optimizer aggregation}\\[-1pt]
{\small 7,500 cluster restarts, exact common scoring, and matched-hardware timing ratios}
\end{center}

\textbf{What this experiment does.}
This report performs no new circuit optimization. It validates and combines
the complete cluster outputs of Experiments 06--08. The inputs contain three
target states, five representations, five circuit depths
$n\in\{4,8,12,16,20\}$, and 100 restarts per target/representation/depth:
\[
3\times5\times5\times100=7{,}500
\]
separately timed optimizer-run records. At fixed target, depth, and restart
index, the five methods intentionally share the same initial circuit; these
five observations are paired, not statistically independent. The aggregation
stops rather than silently using a partial result if any restart or method/depth
pair is absent.

\textbf{Targets.}
Every circuit starts in qubit ground state and oscillator vacuum. The targets
are
\[
\ket{\tau_6}=\frac{\ket9+\ket{10}}{\sqrt2},\qquad
\ket{\tau_8}=\frac{\ket{-2}+\ket2}{\sqrt{2+2e^{-8}}},
\]
and a finite hexagonal GKP logical-zero state. For the latter, define
$\alpha=\sqrt{\pi/\sqrt3}$, $\eta=e^{2\pi i/3}\alpha$, retain integer pairs
$\max(|m|,|l|,|m-l|)\le14$ with even $m$, let $k=m/2$ and
$d_{ml}=m\alpha+l\eta$, and normalize
\[
\ket{\tau_7}\propto\sum_{m,l}
e^{-i\pi kl-0.2^2|d_{ml}|^2}\ket{d_{ml}}.
\]
Here $\ket z$ denotes an oscillator coherent state centered at complex
phase-space position $z$.

\textbf{Five representations.}
Full coherent paths keeps all $2^n$ analytic paths and has no oscillator
cutoff. Analytic Fock carries two 48-entry photon-number vectors and constructs
the projected infinite displacement using
\[
D(a)\ket0=\ket a,\qquad
D(a)\ket r=(\hat a^\dagger-a^*)D(a)\ket{r-1}/\sqrt r.
\]
Dense numerical Fock instead truncates $\hat a$ first and materializes
$D_{48}^{\rm num}(a)=\exp(a\hat a_{48}^\dagger-a^*\hat a_{48})$ and each
$96\times96$ joint layer unitary. Continuous anchors keeps 32 movable coherent
packets per qubit branch, while squeezed packets adds one fitted complex
squeeze per packet. Both packet methods seed from residual pivots, move centers
with inner Adam updates, and globally refit coefficients by
\[
x=(B^\dagger B+10^{-9}I)^{-1}B^\dagger v.
\]
Their discrete selection uses a straight-through gradient approximation.

\textbf{Exact common score and best-of-100 result.}
Write the final joint state as
$\ket\Psi=\sum_{q=0}^1\ket q\ket{\psi_q}$. For target $t$, method $m$,
depth $n$, and restart $r$, the source experiments rescore the saved physical
circuit with
\[
F_{t,m,n,r}=\sum_{q=0}^1|\langle\tau_t|\psi_{q,m,n,r}\rangle|^2.
\]
The qubit is traced out. This score uses analytic coherent-state overlaps and
does not use the Fock or packet approximation. ``Exact'' means exact for the
ideal stated circuit and target up to floating-point arithmetic; hardware noise
is absent. The reported winner and error are
\[
F^*_{t,m,n}=\max_{0\le r<100}F_{t,m,n,r},\qquad
\epsilon^*_{t,m,n}=1-F^*_{t,m,n}.
\]

\IfFileExists{report_results.tex}{\input{report_results.tex}}{
\textbf{Generated results are absent. Run the Experiment 09 workflow first.}}

\textbf{Best quality and restart reliability.}
The heatmap reports $-\log_{10}(1-F^*)$, the number of decimal accuracy digits
in the best observed exact fidelity. A best-of-100 point alone can hide an
unreliable optimizer. For threshold $\tau$, restart reliability is
\[
\widehat p_\tau=\frac1{100}\sum_{r=0}^{99}
\mathbf1[F_{t,m,n,r}\ge\tau].
\]
The second figure uses $\tau=0.99$ and 95\% Wilson binomial intervals.

\begin{figure}[H]
\centering
\includegraphics[width=\textwidth]{figs/best_infidelity_heatmap.png}
\caption{Best exact state-preparation accuracy among 100 restarts. Values are
$-\log_{10}(1-F^*)$; larger is better.}
\end{figure}

\begin{figure}[H]
\centering
\includegraphics[width=\textwidth]{figs/exact_infidelity_vs_layers.png}
\caption{Best exact infidelity. The vertical axis is logarithmic and lower is
better. Every point is selected only by the common exact rescore.}
\end{figure}

\begin{figure}[H]
\centering
\includegraphics[width=\textwidth]{figs/restart_success_fidelity_0p99.png}
\caption{Fraction of the 100 restarts reaching exact fidelity at least 0.99.
Error bars are 95\% Wilson intervals.}
\end{figure}

\textbf{Computer search time.}
For one target, representation, and depth, raw search cost is
\[
C_{t,m,n}=\sum_{r=0}^{99}T_{t,m,n,r},
\]
where $T$ includes optimization and exact rescoring. This is a sum of process
wall times, not Slurm queue time or parallel-array calendar time. Different
targets can run on different CPU models, so raw cross-target times are not
treated as hardware-controlled comparisons.

Within one target, array restart $r$ runs all 25 method/depth cases on one
allocation. For source blocks whose recorded hardware signature is internally
matched, the paired speedup is
\[
S_{t,m,n,r}=\frac{T_{t,{\rm full},n,r}}{T_{t,m,n,r}}.
\]
The plot reports its median and interquartile range. $S>1$ means the method is
faster than full coherent paths on the same target, depth, restart index, and
hardware block.

\begin{figure}[H]
\centering
\includegraphics[width=\textwidth]{figs/total_search_time_vs_layers.png}
\caption{Raw sum of 100 process wall times. Panels keep the targets separate;
the vertical axis is logarithmic.}
\end{figure}

\begin{figure}[H]
\centering
\includegraphics[width=\textwidth]{figs/paired_speedup_vs_full_coherent.png}
\caption{Matched-hardware optimizer speedup relative to full coherent paths.
Lines are medians; bands are interquartile ranges.}
\end{figure}

\textbf{Physical circuit duration.}
Computer search time and hardware execution time are different quantities.
For the winning sequence with ECD amplitudes $\beta_j$, the Coherax estimate is
\[
T_C=\sum_j\left[24\,{\rm ns}+\max\left(
\frac{|\beta_j|}{(2\pi\,50\,{\rm kHz})20},48\,{\rm ns}\right)\right].
\]
The separately labelled Eickbusch-hardware quantity is
\[
T_{E,{\rm LB}}=\sum_j\max\left(
\frac{|\beta_j|}{(2\pi\,33\,{\rm kHz})30},224\,{\rm ns}\right).
\]
The latter is a lower bound, not a compiled waveform duration.

\begin{figure}[H]
\centering
\includegraphics[width=\textwidth]{figs/coherax_gate_time_vs_layers.png}
\caption{Coherax physical gate-time estimate for each best-of-100 circuit.}
\end{figure}

\begin{figure}[H]
\centering
\includegraphics[width=\textwidth]{figs/eickbusch_gate_time_bound_vs_layers.png}
\caption{Eickbusch-hardware lower bound for each best-of-100 circuit.}
\end{figure}

\textbf{Limits of the aggregate.}
These are best-observed random-search outcomes, not global-optimum proofs.
Best-of-100 quality and success probability answer different questions and
must be read together. The geometric mean error in the method table is a
descriptive summary across heterogeneous target states, not a new physical
fidelity. Fock methods retain finite-cutoff effects during training; packet
methods use heuristic projection gradients. Exact final rescoring removes
those approximations from winner evaluation, but cannot remove optimizer bias.

\end{document}
